\documentclass[11pt,a4paper]{article}

\usepackage[utf8]{inputenc}
\usepackage[T1]{fontenc}
\usepackage{XCharter}
\usepackage[brazil]{babel}
\usepackage[margin=2cm]{geometry}
\usepackage{amsmath, amssymb}
\usepackage[xcharter,bigdelims,vvarbb]{newtxmath}
% newtxmath's operators font requests the fontspec family `XCharter(0)' in OT1;
% without these substitutions math digits and operator names silently fall back
% to Computer Modern (LaTeX Font Warning: Font shape `OT1/XCharter(0)/m/n' undefined).
\DeclareFontFamily{OT1}{XCharter(0)}{}
\DeclareFontShape{OT1}{XCharter(0)}{m}{n}{<->ssub * XCharter-TLF/m/n}{}
\DeclareFontShape{OT1}{XCharter(0)}{m}{it}{<->ssub * XCharter-TLF/m/it}{}
\DeclareFontShape{OT1}{XCharter(0)}{b}{n}{<->ssub * XCharter-TLF/b/n}{}
\DeclareFontShape{OT1}{XCharter(0)}{bx}{n}{<->ssub * XCharter-TLF/b/n}{}
\usepackage{enumerate}
\usepackage{array}
\usepackage{tikz}
\usetikzlibrary{positioning}

\pagestyle{empty}
\setlength{\parindent}{0pt}
\setlength{\parskip}{0.5em}

\renewcommand{\arraystretch}{1.4}

\begin{document}

%% =============================================================
%%  Header
%% =============================================================
{\Large\bfseries Aprendizado Profundo} \hfill \textbf{Prova G2, Unidades 1 a 3}\\
Prof. Lucas S. Kupssinskü \hfill Data: \rule{3cm}{0.4pt}\\[0.4em]
Nome: \rule{12cm}{0.4pt}

\rule{\linewidth}{0.8pt}

%% =============================================================
%%  Response grid
%% =============================================================
\textbf{Quadro de respostas (questões objetivas):}\\[0.3em]
\begin{center}
\begin{tabular}{|c|c|c|c|c|c|c|c|c|}
\hline
\textbf{Questão} & \textbf{a} & \textbf{b} & \textbf{c} & \textbf{d} & \textbf{e} & \textbf{f} & \textbf{g} & \textbf{h} \\\hline
1 & & & & & & & & \\\hline
2 & & & & & & & & \\\hline
3 & & & & & & & & \\\hline
4 & & & & & & & & \\\hline
5 & & & & & & & & \\\hline
6 & & & & & & & & \\\hline
7 & & & & & & & & \\\hline
8 & & & & & & & & \\\hline
\end{tabular}
\end{center}

\rule{\linewidth}{0.4pt}

%% =============================================================
%%  Objective questions  (8 x 1,0 = 8,0 pts)
%% =============================================================
\section*{Questões Objetivas (8{,}0 pts)}
\textit{Cada questão possui apenas uma alternativa correta. Marque a letra escolhida no quadro de respostas.}

\begin{enumerate}[1)]

%% ---------- Q1 (U1): OBJ-J plot interpretation — learning-rate curves ----------
\item (1{,}0) Você treina o mesmo MLP quatro vezes, alterando apenas a taxa de aprendizado $\eta$ entre os treinos. As curvas A, B, C e D abaixo mostram a evolução da função de custo $\mathcal{L}$ ao longo das iterações em cada treino (todas partem da mesma inicialização). Assinale a alternativa que associa corretamente cada curva ao regime da taxa de aprendizado.

\begin{center}
\begin{tikzpicture}[>=stealth, every node/.style={font=\small}, xscale=1.05, yscale=0.8]
    \draw[->] (0,0) -- (6.4,0) node[below left] {iterações};
    \draw[->] (0,0) -- (0,3.6) node[above] {$\mathcal{L}$};
    % A: diverges
    \draw[thick] plot[smooth] coordinates {(0,1.8) (0.8,1.55) (1.6,1.9) (2.6,2.5) (3.6,3.3)}
        node[right] {A};
    % B: fast drop then oscillates around plateau
    \draw[thick] plot[smooth] coordinates {(0,1.8) (0.5,1.15) (1.0,0.85) (1.5,1.15) (2.0,0.8) (2.5,1.1) (3.0,0.85) (3.5,1.1) (4.0,0.85) (4.5,1.05) (5.0,0.9) (5.5,1.05)}
        node[right] {B};
    % C: smooth convergence to low value
    \draw[thick] plot[smooth] coordinates {(0,1.8) (0.7,0.95) (1.5,0.5) (2.5,0.3) (3.5,0.22) (4.5,0.19) (5.5,0.18)}
        node[right] {C};
    % D: very slow, almost linear decline
    \draw[thick] plot[smooth] coordinates {(0,1.8) (1.1,1.72) (2.2,1.64) (3.3,1.56) (4.4,1.48) (5.5,1.4)}
        node[right] {D};
\end{tikzpicture}
\end{center}

\begin{enumerate}[a)]
    \item A: $\eta$ alta; \quad B: $\eta$ muito alta (divergência); \quad C: $\eta$ adequada; \quad D: $\eta$ muito baixa.
    \item A: $\eta$ muito alta (divergência); \quad B: $\eta$ adequada; \quad C: $\eta$ alta; \quad D: $\eta$ muito baixa.
    \item A: $\eta$ muito baixa; \quad B: $\eta$ alta; \quad C: $\eta$ adequada; \quad D: $\eta$ muito alta (divergência).
    \item A: $\eta$ muito alta (divergência); \quad B: $\eta$ alta; \quad C: $\eta$ muito baixa; \quad D: $\eta$ adequada.
    \item A: $\eta$ muito alta (divergência); \quad B: $\eta$ alta; \quad C: $\eta$ adequada; \quad D: $\eta$ muito baixa.
    \item A: $\eta$ alta; \quad B: $\eta$ muito baixa; \quad C: $\eta$ adequada; \quad D: $\eta$ muito alta (divergência).
    \item A: $\eta$ muito alta (divergência); \quad B: $\eta$ muito baixa; \quad C: $\eta$ alta; \quad D: $\eta$ adequada.
    \item A: $\eta$ adequada; \quad B: $\eta$ alta; \quad C: $\eta$ muito alta (divergência); \quad D: $\eta$ muito baixa.
\end{enumerate}

%% ---------- Q2 (U1): perceptron logic gate with concrete weights ----------
\item (1{,}0) Considere um Perceptron de Rosenblatt com duas entradas $X_1, X_2 \in \{-1, +1\}$ e saída $\hat{y} = \mathrm{sinal}(\theta_0 + \theta_1 X_1 + \theta_2 X_2)$, em que $\mathrm{sinal}(z) = +1$ se $z \geq 0$ e $-1$ caso contrário. Codificando \emph{verdadeiro} como $+1$ e \emph{falso} como $-1$, qual configuração de pesos $(\theta_0, \theta_1, \theta_2)$ implementa a porta lógica \textbf{OR}, isto é, $\hat{y} = +1$ se e somente se $X_1 = +1$ ou $X_2 = +1$?

\begin{enumerate}[a)]
    \item $(\theta_0, \theta_1, \theta_2) = (-3,\ 1,\ 1)$
    \item $(\theta_0, \theta_1, \theta_2) = (-1,\ -1,\ -1)$
    \item $(\theta_0, \theta_1, \theta_2) = (-1,\ 1,\ 1)$
    \item $(\theta_0, \theta_1, \theta_2) = (1,\ 1,\ 1)$
    \item $(\theta_0, \theta_1, \theta_2) = (1,\ 1,\ -1)$
    \item $(\theta_0, \theta_1, \theta_2) = (1,\ -1,\ -1)$
    \item $(\theta_0, \theta_1, \theta_2) = (-1,\ 1,\ -1)$
    \item $(\theta_0, \theta_1, \theta_2) = (3,\ 1,\ 1)$
\end{enumerate}

%% ---------- Q3 (U1): numeric — two momentum iterations ----------
\item (1{,}0) Um único parâmetro $\theta$ é otimizado por descida de gradiente com \textit{momentum}, na formulação vista em aula:
\[
v_{t+1} = \beta\,v_{t} + (1-\beta)\,\nabla_{\theta}\mathcal{L}(\theta_{t}), \qquad \theta_{t+1} = \theta_{t} - \eta\,v_{t+1},
\]
com $\eta = 0{,}1$, $\beta = 0{,}5$, $v_0 = 0$ e $\theta_0 = 1$. Os gradientes observados nas duas primeiras iterações são $\nabla_{\theta}\mathcal{L}(\theta_0) = 4$ e $\nabla_{\theta}\mathcal{L}(\theta_1) = 2$. Qual o valor de $\theta_2$ após \textbf{duas} atualizações?

\begin{enumerate}[a)]
    \item $0{,}2$
    \item $0{,}4$
    \item $0{,}5$
    \item $0{,}6$
    \item $0{,}7$
    \item $0{,}8$
    \item $1{,}0$
    \item $1{,}4$
\end{enumerate}

%% ---------- Q4 (U2): OBJ-E lacunas — vision tasks ----------
\item (1{,}0) Assinale a alternativa que preenche corretamente as lacunas do parágrafo abaixo.

\begin{quote}
``Para medir a qualidade de uma \textit{bounding box} predita em relação ao \textit{ground truth}, utiliza-se tipicamente a métrica \textbf{(1)}. A tarefa que atribui um rótulo de classe a cada \textit{pixel} da imagem, sem distinguir objetos diferentes de uma mesma classe, chama-se segmentação \textbf{(2)}. Ao adaptar uma VGG-16 pré-treinada no ImageNet para um novo problema com poucos dados, a prática usual de \textit{transfer learning} é congelar \textbf{(3)} e treinar o restante da rede.''
\end{quote}

\begin{enumerate}[a)]
    \item (1) entropia cruzada; (2) semântica; (3) o \textit{backbone} convolucional
    \item (1) IoU; (2) de instâncias; (3) o \textit{backbone} convolucional
    \item (1) entropia cruzada; (2) de instâncias; (3) a cabeça densa
    \item (1) IoU; (2) semântica; (3) a cabeça densa
    \item (1) entropia cruzada; (2) semântica; (3) a cabeça densa
    \item (1) IoU; (2) de instâncias; (3) a cabeça densa
    \item (1) IoU; (2) semântica; (3) o \textit{backbone} convolucional
    \item (1) entropia cruzada; (2) de instâncias; (3) o \textit{backbone} convolucional
\end{enumerate}

%% ---------- Q5 (U2): OBJ-F assertion-reason — degradation problem ----------
\item (1{,}0) Considere as asserções abaixo e a relação proposta entre elas.

\textbf{I.} Ao aumentar muito a profundidade de uma rede convolucional \textbf{sem} conexões residuais (por exemplo, de 20 para 56 camadas), a rede mais profunda pode apresentar desempenho \textbf{pior} que a mais rasa; esse fenômeno, conhecido como problema da degradação, não é um caso de \textit{overfitting}.

\begin{center}\textbf{PORQUE}\end{center}

\textbf{II.} Na degradação, o erro elevado da rede mais profunda aparece já no conjunto de \textbf{treino}, enquanto o \textit{overfitting} se caracteriza por erro de treino baixo acompanhado de erro de validação alto.

Assinale a alternativa correta:
\begin{enumerate}[a)]
    \item As asserções I e II são verdadeiras, e a II é uma justificativa correta da I.
    \item As asserções I e II são verdadeiras, mas a II não é uma justificativa correta da I.
    \item As asserções I e II são verdadeiras, e a I é uma justificativa correta da II, mas não o contrário.
    \item A asserção I é verdadeira, e a II é falsa.
    \item A asserção I é falsa, e a II é verdadeira.
    \item As asserções I e II são falsas.
    \item A asserção I é verdadeira somente quando a rede usa \textit{batch normalization}, e a II é verdadeira.
    \item A asserção I é falsa somente para redes treinadas com SGD sem \textit{momentum}, e a II é falsa.
\end{enumerate}

%% ---------- Q6 (U3): OBJ-C scenario diagnostic — exploding gradients ----------
\item (1{,}0) Você treina uma RNN \textit{vanilla} (ativação $\tanh$ nas conexões recorrentes) para modelagem de linguagem em sequências longas. A curva da função de custo cai durante algumas centenas de iterações, mas de tempos em tempos \textbf{salta abruptamente} para valores muito altos e, em uma dessas ocasiões, passa a registrar \texttt{NaN}. Ao imprimir a norma do gradiente $\|\mathbf{g}\|$ a cada iteração, você observa que ela ocasionalmente atinge valores na ordem de $10^{6}$. Qual a explicação mais consistente \textbf{e} a melhor mitigação?

\begin{enumerate}[a)]
    \item Os gradientes se dissipam ao longo dos passos de tempo; substituir a RNN por uma LSTM elimina os saltos, pois os portões impedem qualquer crescimento da norma do gradiente.
    \item A variância do gradiente está alta por causa do tamanho do \textit{batch}; reduzir o \textit{batch size} para $1$ torna as atualizações mais estáveis e elimina os saltos e o \texttt{NaN} observados.
    \item O gradiente explode ocasionalmente; aplicar \textit{clipping} por valor em cada componente é preferível ao \textit{clipping} por norma, pois somente o primeiro preserva a direção do vetor gradiente.
    \item O modelo está em \textit{overfitting} nas sequências longas; adicionar \textit{dropout} nas conexões recorrentes reduz a variância do gradiente e impede que a norma atinja valores extremos.
    \item A saturação da $\tanh$ é a causa dos saltos; trocá-la por ReLU nas conexões recorrentes limita o crescimento das ativações e, portanto, o crescimento da norma do gradiente.
    \item O gradiente explode: a retropropagação no tempo multiplica fatores $\theta_{hh}\,(1-\tanh^2(z_k))$ que podem ter módulo maior que $1$; aplicar \textit{clipping} por norma ($\mathbf{g} \leftarrow c\,\mathbf{g}/\|\mathbf{g}\|$ quando $\|\mathbf{g}\| > c$) limita o passo preservando a direção.
    \item O \texttt{NaN} indica erro numérico interno do \textit{framework} de autodiferenciação; reexecutar o treinamento com outra semente aleatória, sem alterar hiperparâmetros, é suficiente para eliminar o problema.
    \item O treinamento com \textit{teacher forcing} acumula erros de exposição; alimentar o modelo com as próprias predições durante o treino estabiliza a norma do gradiente e evita os saltos.
\end{enumerate}

%% ---------- Q7 (U3): OBJ-H matching — decoding strategies ----------
\pagebreak
\item (1{,}0) Associe cada estratégia de decodificação de um modelo de linguagem (coluna da esquerda) à descrição correspondente (coluna da direita).

\begin{center}
\begin{tabular}{p{0.42\linewidth} p{0.5\linewidth}}
\textbf{1.} \textit{Greedy search} & \textbf{(i)} mantém, a cada passo, as $b$ hipóteses de maior probabilidade conjunta e retorna a melhor ao final. \\
\textbf{2.} \textit{Beam search} (largura $b$) & \textbf{(ii)} escolhe deterministicamente, a cada passo, o \textit{token} de maior probabilidade condicional. \\
\textbf{3.} Amostragem com temperatura $\tau \to 0^{+}$ & \textbf{(iii)} torna-se equivalente à escolha determinística do \textit{token} mais provável a cada passo. \\
\textbf{4.} Amostragem Top-P (\textit{nucleus}) & \textbf{(iv)} amostra apenas do menor conjunto de \textit{tokens} cuja probabilidade acumulada atinge um limiar $p$. \\
\end{tabular}
\end{center}

\begin{enumerate}[a)]
    \item 1-(ii); \quad 2-(i); \quad 3-(iv); \quad 4-(iii)
    \item 1-(iii); \quad 2-(i); \quad 3-(ii); \quad 4-(iv)
    \item 1-(ii); \quad 2-(i); \quad 3-(iii); \quad 4-(iv)
    \item 1-(i); \quad 2-(ii); \quad 3-(iii); \quad 4-(iv)
    \item 1-(ii); \quad 2-(iv); \quad 3-(iii); \quad 4-(i)
    \item 1-(i); \quad 2-(ii); \quad 3-(iv); \quad 4-(iii)
    \item 1-(iii); \quad 2-(iv); \quad 3-(ii); \quad 4-(i)
    \item 1-(iv); \quad 2-(i); \quad 3-(iii); \quad 4-(ii)
\end{enumerate}

%% ---------- Q8 (U3): numeric — dot-product attention weight ----------
\item (1{,}0) Em um modelo \textit{seq2seq} com atenção de produto escalar, o estado do decodificador no passo $t$ é $\mathbf{h}_t = (1;\ 1)$ e os estados do codificador são $\mathbf{s}_1 = (2;\ 2)$, $\mathbf{s}_2 = (1;\ 1)$ e $\mathbf{s}_3 = (0;\ -1)$. Os \textit{scores} são $\mathrm{score}(\mathbf{h}_t, \mathbf{s}_k) = \mathbf{h}_t^{\top} \mathbf{s}_k$ e os pesos de atenção são
\[
a_k^{(t)} = \frac{\exp(\mathrm{score}(\mathbf{h}_t, \mathbf{s}_k))}{\sum_{i=1}^{3}\exp(\mathrm{score}(\mathbf{h}_t, \mathbf{s}_i))}.
\]
Qual o valor aproximado de $a_1^{(t)}$?

\begin{enumerate}[a)]
    \item $0{,}006$
    \item $0{,}118$
    \item $0{,}333$
    \item $0{,}571$
    \item $0{,}800$
    \item $0{,}876$
    \item $0{,}881$
    \item $0{,}982$
\end{enumerate}

\end{enumerate}

\rule{\linewidth}{0.4pt}

%% =============================================================
%%  Dissertative questions  (2 x 1,0 = 2,0 pts)
%% =============================================================
\pagebreak
\section*{Questões Dissertativas (2{,}0 pts)}

\begin{enumerate}[1)]
\setcounter{enumi}{8}

%% ---------- Q9 (U3): DISS-A — LSTM single step by hand ----------
\item (1{,}0) Considere uma célula LSTM \textbf{escalar} (todas as dimensões iguais a $1$), com as equações vistas em aula:
\[
f_t = \sigma(\boldsymbol{\theta}_f\,[h_{t-1}, x_t] + b_f), \qquad
i_t = \sigma(\boldsymbol{\theta}_i\,[h_{t-1}, x_t] + b_i), \qquad
\tilde{c}_t = \tanh(\boldsymbol{\theta}_c\,[h_{t-1}, x_t] + b_c),
\]
\[
c_t = f_t \odot c_{t-1} + i_t \odot \tilde{c}_t, \qquad
o_t = \sigma(\boldsymbol{\theta}_o\,[h_{t-1}, x_t] + b_o), \qquad
h_t = o_t \odot \tanh(c_t).
\]
O estado inicial é $c_0 = -0{,}2$ e $h_0 = 0{,}4$, e a entrada no passo $1$ é $x_1 = 0{,}5$, de modo que $[h_0, x_1] = [0{,}4,\ 0{,}5]$. Os parâmetros são:
\[
\boldsymbol{\theta}_f = [0{,}5,\ -0{,}4],\ b_f = 0{,}2; \qquad
\boldsymbol{\theta}_i = [-0{,}3,\ 0{,}6],\ b_i = -0{,}1; \qquad
\]
\[
\boldsymbol{\theta}_c = [0{,}2,\ 0{,}8],\ b_c = 0; \qquad
\boldsymbol{\theta}_o = [0{,}7,\ 0{,}1],\ b_o = -0{,}3.
\]
Pode arredondar para $4$ casas decimais.

\begin{enumerate}[a)]
    \item (0{,}5) Compute $f_1$, $i_1$, $\tilde{c}_1$ e o novo estado de célula $c_1$.
    \vspace{3.4cm}

    \item (0{,}5) Compute $o_1$ e $h_1$. Em seguida, explique em uma frase, a partir de $\dfrac{\partial c_t}{\partial c_{t-1}} = f_t$, por que a LSTM mitiga a dissipação do gradiente em sequências longas.
    \vspace{3.4cm}
\end{enumerate}

%% ---------- Q10 (U1+U2): DISS-B — MLP vs CNN for images ----------
\item (1{,}0) Um colega propõe usar um MLP para classificar imagens RGB de dimensão $64 \times 64 \times 3$, achatando a imagem em um vetor de entrada e usando uma primeira camada oculta totalmente conectada com $256$ unidades.

\begin{enumerate}[a)]
    \item (0{,}5) Compute o número de parâmetros treináveis (pesos e vieses) dessa primeira camada totalmente conectada e compare-o com o número de parâmetros de uma camada convolucional $3 \times 3$ com $32$ filtros (com viés), \textit{stride} $1$ e \textit{padding} $1$, aplicada à mesma imagem. Mostre o cálculo dos dois valores.
    \vspace{3.4cm}

    \item (0{,}5) Além da mudança no número de parâmetros, aponte \textbf{duas} propriedades estruturais da convolução que a tornam adequada para imagens, explicando em uma ou duas frases por que cada uma delas ajuda nessa tarefa.
    \vspace{3.4cm}
\end{enumerate}

\end{enumerate}

\end{document}
